
\begin{obeassessment}
\textbf{Kuis Bab 13}
1. Apa perbedaan antara persoalan \textit{NP-complete} dan persoalan yang dapat
   diselesaikan dalam waktu polinomial?
2. Jelaskan langkah-langkah transformasi dari kunci privat ke kunci publik pada
   sistem Merkle-Hellman.
3. Mengapa penemuan Adi Shamir membuat algoritma knapsack tidak lagi aman?
4. Jelaskan mengapa algoritma \textit{greedy} selalu berhasil pada barisan
   \textit{superincreasing}. Jalankan algoritma itu pada barisan
   $\{1, 5, 6, 11, 14, 20\}$ untuk $M = 32$. Apakah hasilnya sama dengan blok
   \code{111001}? Apa yang ditunjukkan keadaan itu?
5. Diketahui $W_{priv} = \{2, 3, 6, 13, 27, 52\}$, $m = 105$, dan $n = 31$.
   Hitunglah $W_{pub}$ dan $n^{-1} \bmod 105$, lalu enkripsikan blok
   \code{110101} dan periksa hasil dekripsinya.
6. Mengapa $C = \sum_i b_i w'_i$ tidak perlu diambil modulo $m$, padahal
   $w'_i$ sendiri dihitung modulo $m$? Kaitkan jawaban Anda dengan syarat
   $m > \sum_i w_i$.
7. Sumber menganjurkan kunci berisi 250 elemen berukuran 200 sampai 400 bit
   dan modulus 100 sampai 200 bit. Mengapa kedua anjuran itu tidak konsisten?
   Berapa panjang modulus yang diperlukan?
\end{obeassessment}
\begin{competencychecklist}
\begin{tabularx}{\linewidth}{|X|c|}
\hline
Indikator Kompetensi & Tercapai \\
\hline
Saya dapat menjelaskan konsep \textit{Knapsack Problem} & [ ] \\
\hline
Saya dapat membedakan barisan \textit{superincreasing} dan normal & [ ] \\
\hline
Saya dapat memeriksa apakah sebuah barisan bersifat \textit{superincreasing} & [ ] \\
\hline
Saya dapat memecahkan \textit{superincreasing knapsack} dengan algoritma \textit{greedy} & [ ] \\
\hline
Saya dapat menghitung kunci publik dari kunci privat Merkle-Hellman & [ ] \\
\hline
Saya dapat menjelaskan mengapa $m > \sum_i w_i$ dan $\text{PBB}(n,m)=1$ wajib dipenuhi & [ ] \\
\hline
Saya dapat menghitung invers modulo $n^{-1} \bmod m$ & [ ] \\
\hline
Saya dapat menghitung enkripsi sistem Merkle-Hellman & [ ] \\
\hline
Saya dapat melakukan dekripsi dengan bantuan invers modulo & [ ] \\
\hline
Saya dapat menjelaskan mengapa dekripsi Merkle-Hellman berhasil & [ ] \\
\hline
Saya dapat menganalisis penyebab kegagalan keamanan algoritma knapsack & [ ] \\
\hline
\end{tabularx}
\end{competencychecklist}
